\section{Congestion-Aware Deferral}\label{sec:congestion}

If contention is the main cause of harm, the remedy should target it. Our simulator already makes
every deferred job reserve its target slot, so later jobs see the reduced capacity. That alone does not
stop deferred jobs from filling the cheapest hours and leaving no room for work that arrives during them.
We test two minimal changes to point-forecast greedy.

\textbf{P8, headroom-reserved deferral.} A \emph{deferred} reservation may use at most $(1-h)$ of any
slot's capacity; running now may use all of it. With $h{=}0.25$, a quarter of every future hour stays
available to arrivals. The carbon model is unchanged. P8's capacity reservation is a practical,
online approximation of the joint assignment that a global optimum would make: it makes deferral
decisions contention-aware without global foresight.

\textbf{P9, jittered deferral.} Instead of the argmin, pick uniformly among the three best slots that
are forecast no worse than running now (P9-r re-plans but keeps its slot while it stays in the top three).

\textbf{Results.} At 70\% utilisation, P8 lowers the harm rate from
\res{V2U}{ALL-S12-utiltarget0.7}{P3}{harmrate}\% to \res{V2U}{ALL-S12-utiltarget0.7}{P8}{harmrate}\%
(mean of six regions) for \res{V2X}{ALL-u0.7}{P8}{savinglossvsPthree} points of saving. Harm is lower
in \res{V2W}{u0.7}{P8vsP3}{harmratewins} of \res{V2W}{u0.7}{P8vsP3}{harmraten} paired region-weeks
(Wilcoxon $p{<}0.05$ in \res{V2W}{u0.7}{P8vsP3}{harmratesig} of 6 regions). On GB-op it saves
\res{V2U}{GBop-S12-utiltarget0.7}{P8}{savingpct}\% (P3: \res{V2U}{GBop-S12-utiltarget0.7}{P3}{savingpct}\%)
while cutting harm by \res{V2X}{GBop-u0.7}{P8}{harmcutvsPthree}\%. It also shortens the tail: p95
slowdown \res{V2U}{ALL-S12-utiltarget0.7}{P8}{slowdownp95} vs.\ \res{V2U}{ALL-S12-utiltarget0.7}{P3}{slowdownp95},
and non-deferrable jobs wait \res{V2U}{ALL-S12-utiltarget0.7}{P8}{waitminmeannondeferrable} instead of
\res{V2U}{ALL-S12-utiltarget0.7}{P3}{waitminmeannondeferrable} minutes on average. The mechanism is visible
in the split: P8's \emph{voluntary} harm is not lower
(\res{V2U}{ALL-S12-utiltarget0.7}{P8}{harmratevoluntary}\% vs.\ \res{V2U}{ALL-S12-utiltarget0.7}{P3}{harmratevoluntary}\%);
what disappears is forced harm on jobs that found the cluster full.
In Fig.~\ref{fig:dial}, P8 lies above the conformal dial, and it is non-dominated by any dial point in
\res{V2F}{ALL}{P8}{nondominated} of \res{V2F}{ALL}{P8}{nregions} regions: there, no uncertainty-aware
setting achieves both lower harm and higher saving.

\textbf{Limits.} P8 does not reach the demanding target of P4's harm at P3's savings: it meets both in
only \res{V2F}{u0.7}{P8}{hitstarget} region at 70\% and in none at 50\% or 90\%. Emissions under P8 are
slightly higher than under P3 in most weeks (significantly so in
\res{V2W}{u0.7}{P8vsP3}{emissionskgsig} of 6 regions, but not on GB-op). At 90\% utilisation the saving
drops to \res{V2U}{ALL-S12-utiltarget0.9}{P8}{savingpct}\% against
\res{V2U}{ALL-S12-utiltarget0.9}{P3}{savingpct}\%. A larger headroom ($h{=}0.5$) over-restricts deferral
(saving \res{V2U}{ALL-S12-utiltarget0.7}{P8-h0.5}{savingpct}\%). Jitter barely helps: P9-r's harm rate is
\res{V2U}{ALL-S12-utiltarget0.7}{P9-r}{harmrate}\% against P3's
\res{V2U}{ALL-S12-utiltarget0.7}{P3}{harmrate}\%, because the three best slots are usually adjacent hours
of the same trough. Herding onto a single argmin is therefore not the whole story. What matters is how
much of the trough deferred work may occupy.
